%!TEX root = ../main.tex
% APÉNDICE B — Entrevistas a especialistas

\chapter{Entrevistas a especialistas}\label{apx:entrevista}

\begin{guia}
Un apéndice por entrevista: contexto (fecha, medio, duración, objetivo),
datos de la persona entrevistada (formación, cargo, trayectoria, con su
consentimiento) y la transcripción organizada por temas. Las
conclusiones van en el cuerpo (capítulos \ref{cap:marco-teorico} y
\ref{cap:resultados}); acá va la evidencia.
\end{guia}

\section{Entrevista a responsable de salón de eventos}

\subsection{Contexto de la entrevista}

Fecha de realización: 22/07/2026. \completar{Modalidad, duración y
objetivo específico de la entrevista.}

\subsection{Datos de la persona entrevistada}

Dueño del salón, con más de 10 años de experiencia en el sector de
eventos. El salón realiza entre 51 y 100 eventos por año (casamientos,
cumpleaños, fiestas de 15 años y eventos empresariales).

\subsection{Transcripción de la entrevista}

\begin{description}
    \item[Entrevistador] ¿Qué cargo o función desempeña dentro del salón?
    \item[Entrevistado] Dueño.
    \item[Entrevistador] ¿Cuántos años de experiencia posee en el sector
          de eventos?
    \item[Entrevistado] Más de 10 años.
    \item[Entrevistador] ¿Cuántos eventos realiza aproximadamente el
          salón por año?
    \item[Entrevistado] Entre 51 y 100 eventos.
    \item[Entrevistador] ¿Qué tipos de eventos se realizan habitualmente
          en el salón?
    \item[Entrevistado] Casamientos, cumpleaños, fiestas de 15 años y
          eventos empresariales.
    \item[Entrevistador] Describa brevemente el proceso desde la consulta
          hasta la finalización del evento.
    \item[Entrevistado] Llega la consulta por WhatsApp, se envía el
          presupuesto, se coordina una visita y finalmente se firma el
          contrato y se abona la seña.
    \item[Entrevistador] ¿Qué información considera indispensable
          registrar para cada evento?
    \item[Entrevistado] Fecha del evento, cantidad de invitados, pagos y
          servicios contratados.
    \item[Entrevistador] ¿Qué herramientas utiliza actualmente para
          organizar los eventos?
    \item[Entrevistado] WhatsApp y Microsoft Excel.
    \item[Entrevistador] ¿La información de cada evento se encuentra
          centralizada en un único lugar?
    \item[Entrevistado] Parcialmente; una parte está centralizada y otra
          distribuida.
    \item[Entrevistador] ¿Cuáles son las principales dificultades que
          enfrenta durante la organización de un evento?
    \item[Entrevistado] La coordinación de proveedores.
    \item[Entrevistador] ¿Qué problemas aparecen con mayor frecuencia al
          coordinar cliente, salón y proveedores?
    \item[Entrevistado] La coordinación entre los distintos proveedores.
    \item[Entrevistador] ¿Ha tenido inconvenientes por información
          incorrecta o cambios de último momento?
    \item[Entrevistado] Cambios de horarios y contratación de
          proveedores.
    \item[Entrevistador] ¿Qué información le resulta más difícil de
          mantener actualizada?
    \item[Entrevistado] Los pagos.
    \item[Entrevistador] ¿Qué funcionalidades debería tener una
          plataforma para facilitar la gestión de un salón?
    \item[Entrevistado] Calendario, registro de pagos recibidos y pagos a
          proveedores.
    \item[Entrevistador] ¿Le resultaría útil que clientes y proveedores
          consultaran información actualizada según su función?
    \item[Entrevistado] Sí, sería muy útil.
    \item[Entrevistador] ¿Le resultaría útil que la plataforma recomendara
          proveedores según tipo de evento, ubicación, presupuesto y
          disponibilidad?
    \item[Entrevistado] No está seguro.
    \item[Entrevistador] ¿Considera útil una simulación del evento antes
          de realizarlo?
    \item[Entrevistado] Considera que sería útil principalmente para los
          clientes.
\end{description}

\section{Entrevista a profesional de decoración y ambientación}

\subsection{Contexto de la entrevista}

Fecha de realización: 23/07/2026. \completar{Modalidad, duración y
objetivo específico de la entrevista.}

\subsection{Datos de la persona entrevistada}

Gerente de planificación y ejecución, con más de 10 años de experiencia
en el sector de eventos. Trabaja en entre 26 y 50 eventos por año
(casamientos, cumpleaños, fiestas de 15 años, eventos empresariales e
institucionales).

\subsection{Transcripción de la entrevista}

\begin{description}
    \item[Entrevistador] ¿Qué cargo o función desempeña?
    \item[Entrevistado] Gerente de planificación y ejecución.
    \item[Entrevistador] ¿Cuántos años de experiencia posee en el sector
          de eventos?
    \item[Entrevistado] Más de 10 años.
    \item[Entrevistador] ¿Cuántos eventos realiza aproximadamente por
          año?
    \item[Entrevistado] Entre 26 y 50 eventos.
    \item[Entrevistador] ¿En qué tipos de eventos trabaja habitualmente?
    \item[Entrevistado] Casamientos, cumpleaños, fiestas de 15 años,
          eventos empresariales e institucionales.
    \item[Entrevistador] Describa el proceso desde la consulta hasta el
          desmontaje.
    \item[Entrevistado] Se presenta una propuesta inicial que se
          personaliza según el cliente; se adquiere la materia prima, se
          prepara la ambientación y un equipo de tres o cuatro personas
          realiza el armado y desarme.
    \item[Entrevistador] ¿Qué información considera indispensable antes
          de elaborar una propuesta?
    \item[Entrevistado] Toda la información mencionada anteriormente.
    \item[Entrevistador] ¿Qué herramientas utiliza actualmente para
          organizar sus trabajos?
    \item[Entrevistado] WhatsApp, Microsoft Excel, agenda en papel,
          Word/Google Docs y redes sociales.
    \item[Entrevistador] ¿La información de cada evento se encuentra
          centralizada?
    \item[Entrevistado] Sí, toda la información se encuentra
          centralizada.
    \item[Entrevistador] ¿Cuáles son las principales dificultades que
          enfrenta?
    \item[Entrevistado] Generalmente no se presentan dificultades cuando
          existe buena comunicación entre las partes.
    \item[Entrevistador] ¿Qué problemas aparecen con mayor frecuencia al
          coordinar con el salón y otros proveedores?
    \item[Entrevistado] Coordinar correctamente los horarios de armado y
          desarme.
    \item[Entrevistador] ¿Ha tenido inconvenientes por cambios de último
          momento?
    \item[Entrevistado] No; mantienen comunicación constante y adecuada
          coordinación.
    \item[Entrevistador] ¿Qué información le resulta más difícil de
          mantener actualizada?
    \item[Entrevistado] Los precios de trabajos adicionales o
          instalaciones especiales.
    \item[Entrevistador] ¿Con qué frecuencia realiza visitas previas al
          salón?
    \item[Entrevistado] Considera la visita técnica muy importante; un
          plano a escala ayudaría a mejorar la planificación.
    \item[Entrevistador] ¿Qué funcionalidades debería tener una
          plataforma para la gestión de proyectos de decoración?
    \item[Entrevistado] Inventario actualizado y portafolio de imágenes
          con precios de las propuestas.
    \item[Entrevistador] ¿Le resultaría útil consultar información
          actualizada del evento según su función?
    \item[Entrevistado] Sí, sería muy útil.
    \item[Entrevistador] ¿Le resultaría útil que la plataforma
          recomendara proveedores adecuados?
    \item[Entrevistado] Sí, sería muy útil.
    \item[Entrevistador] ¿Considera útil una simulación previa de la
          ambientación?
    \item[Entrevistado] Considera que aportaría mayor seguridad al
          cliente durante la planificación.
\end{description}

\section{Entrevistas pendientes}

\completar{Entrevista a responsable de catering y encuesta al público:
faltan para completar la Fase~2 de la metodología
(capítulo~\ref{cap:metodologia}). Agregar acá en cuanto estén realizadas.}
